On the ring shaped domain
\(
\Omega
=
\{
(\varphi,r)
\mid
0\le \varphi < 2\pi
\wedge
1<r<2
\}
\)
in polar coordinates, consider the partial differential equation
\begin{equation}
\Delta u = \lambda u
\label{40000035:pde}
\end{equation}
with homogeneous boundary conditions.
Convert this equation into two ordinary differential equations
with appropriate boundary conditions and solve at least one of them.

\begin{hinweis}
Note that the ring shape of the domain implies periodic boundary conditions
for the variable $\varphi$.
\end{hinweis}

\begin{loesung}
The Laplace operator in polar coordinates has the form
\[
\Delta
=
\frac1r \frac{\partial}{\partial r}\frac{\partial}{\partial r}
+
\frac{1}{r^2}\frac{\partial^2}{\partial\varphi^2}
=
\frac{\partial^2}{\partial r^2} + \frac1r\frac{\partial}{\partial r}
+
\frac{1}{r^2}\frac{\partial^2}{\partial\varphi^2}.
\]
To separate $\varphi$ and $r$, we write the solution as a product
$u(\varphi,r)=\Phi(\varphi)R(r)$ and substitute this ansatz into the
differential equation~\eqref{40000035:pde}.
This gives
\[
R''(r)\Phi(\varphi)
+
\frac1rR'(r)\Phi(\varphi)
+
\frac1{r^2}R(r)\Phi''(\varphi)
=
\lambda R(r)\Phi(\varphi).
\]
Dividing by $R(r)\Phi(\varphi)$ allows to separate $\varphi$ and $r$
as follows:
\begin{align*}
\frac{R''(r)}{R(r)}
+
\frac{1}{r}\frac{R'(r)}{R(r)}
+
\frac{1}{r^2} \frac{\Phi''(\varphi)}{\Phi(\varphi)}
&=
\lambda.
\intertext{We need to remove the $r$ in the $\Phi$-term by multiplying the
equation by $r^2$:}
r^2 \frac{R''(r)}{R(r)}
+
r \frac{R'(r)}{R(r)}
-
\lambda r^2
&=
-\frac{\Phi''(\varphi)}{\Phi(\varphi)}.
\end{align*}
The left hand side only depends on $r$, the right hand side only on
$\varphi$, so they must both be constant.
As we need periodic solutions for $\Phi(\varphi)$, the constant
must be positive, we write it as $\mu^2$.
Thus we get the two differential equations
\begin{align}
\Phi''(\varphi) &= -\mu^2 \Phi(\varphi)
&
r^2 R''(r) + rR'(r) -\lambda r^2 R(r) &= \mu^2 R(r).
\label{40000035:separated}
\end{align}
The first equation can be solved by trigonometric functions
\[
\Phi(\varphi)
=
\begin{cases}
\cos \mu \varphi \\
\sin \mu \varphi.
\end{cases}
\]
Periodic boundary conditions imply that
$\cos\mu\varphi=\cos(\varphi+2\pi)$
and
$\sin\mu\varphi=\sin(\varphi+2\pi)$
for all $\varphi$.
Using the addition formula for the trigonometric functions, we get
\[
\renewcommand{\arraycolsep}{2pt}
\begin{array}{rcrcl}
\cos\mu\varphi &=& \cos\mu\varphi\cos 2\mu\pi &-& \sin\mu\varphi\sin2\mu\pi
\\
\sin\mu\varphi &=& \sin\mu\varphi\cos 2\mu\pi &+& \cos\mu\varphi\sin2\mu\pi.
\end{array}
\]
In both cases, this is only possible if $\cos2\mu\pi=1$ and
$\sin2\mu\pi=0$.
The second condition imples that $2\mu\pi$ must be an integer
multiple of $\pi$, but if it is an odd multiple of $\pi$, then
$\cos2\mu\pi=-1$ and first condition would be violated.
Consequently, $2\mu\pi$ must be an even multiple of $\pi$ or $\mu$
must be an integer.
Thus the solutions of the first equation are the functions
\[
\cos k\varphi
\qquad\text{and}\qquad
\sin k\varphi
\qquad\text{with $k\in\mathbb{N}$.}
\]
The second differential equation in \eqref{40000035:separated} is 
\[
r^2R''(r)
+
r R'(r)
-
(k^2-\lambda r^2)R(r)
=
0,
\]
this is Bessel's equation, which can be solved using Bessel functions.
\end{loesung}

\begin{bewertung}
Product ansatz ({\bf A}) 1 point,
separation of variables ({\bf S}) 1 point,
two separated differential equations ({\bf D}) 2 point,
solutions of the equations ({\bf L}) 1 point,
restriction of the value of $\mu$ from boundary conditions ({\bf B}) 1 point.
\end{bewertung}

